\appendix
\chapter{Reproduction and Artifact Manifest}

The repository root contains \texttt{REPRODUCING.md}, \texttt{evidence-manifest.json}, the locked dependency graph, Docker configurations, and the versioned corpus. Legacy reports are explicitly excluded. A final accepted manifest must contain the clean Git commit, Docker digests, corpus hash, run identifiers, and hashes of generated thesis tables.

\section{Command Sequence}

\begin{Verbatim}[breaklines=true,fontsize=\small]
uv sync --locked --group dev
uv run ruff check src tests scripts
uv run mypy src/unified_multi_agent_coordination
uv run pytest --cov
uv run unified-coordination-scenarios --no-output
docker compose -f docker-compose.system.yml up --build \
  --abort-on-container-exit --exit-code-from system-tests
docker compose -f docker-compose.distributed.yml up --build \
  --abort-on-container-exit --exit-code-from distributed-system-tests
docker compose -f docker-compose.etcd-distributed.yml up --build \
  --abort-on-container-exit --exit-code-from etcd-system-tests
uv run unified-defense-study --check
uv run unified-analyze-defense-study \
  demo_runs/v0.3/20260713T013624Z-b9fd8b6f83 \
  --corpus corpus/v0.3 \
  --output demo_runs/v0.3/20260713T013624Z-b9fd8b6f83/analysis.json
uv run unified-defense-study-v04 --corpus corpus/v0.4 --check
uv run unified-analyze-defense-study-v04 \
  --run demo_runs/v0.4/20260713T121506Z-a0ecd2bf0e \
  --corpus corpus/v0.4
uv run unified-runtime-ablation-study \
  --output-dir demo_runs/runtime_ablation/<new-run-id>
\end{Verbatim}

\chapter{Predicate Refinement and Claim--Evidence Matrix}

\begin{longtable}{p{0.22\textwidth}p{0.35\textwidth}p{0.33\textwidth}}
\toprule
\textbf{Claim or predicate} & \textbf{Implementation refinement} & \textbf{Evidence required} \\
\midrule
Coverage & Canonical capability ID, full modes and schema coverage & Feasibility and ID-perturbation tests \\
Authority & Local trust order and admission policy & Remote-card and credential tests \\
Constraints & Typed source, JSON pointer, operator and expected value & Missing, type, false and unresolved cases \\
Verifiability & Task, requirement or capability validation contract & Validatorless and nested-schema tests \\
Dependency safety & Completed prerequisites and declared artifacts only & Failure/timeout/unknown branch tests \\
Traceability & Persisted session, plan, task and attempt events & Ordering, restart and isolation tests \\
Recovery & Lease, fencing, idempotency and unknown-side-effect policy & PostgreSQL crash-fault harness \\
Distributed authority & etcd/Raft state, learner-first membership, odd voter target and quorum fail-closed behavior & Three-voter report plus planned partition and resize matrix \\
RQ4 comparison & Author-labelled, label-hidden v0.3 and v0.4 studies & Two immutable 540-output matrices and reproducible digests \\
Runtime mechanisms & Secure-default policy surfaces and explicit negative controls & Separate 350-observation runtime-only bundle \\
\bottomrule
\end{longtable}

\chapter{Tribunal Questions and Defensible Boundaries}

\begin{description}
  \item[What is novel?] The protocol-adaptable integration and operational refinement of a linguistic proposal boundary, symbolic authorization/refusal, durable evidence, admission policy, and measured recovery. Individual mechanisms are prior art.
  \item[Why use an LLM?] Natural-language interpretation and candidate decomposition vary semantically; a structured rule oracle assumes that work is already done. The LLM remains non-authoritative.
  \item[What is guaranteed?] The tested implementation enforces declared predicates and ordering under its trusted-code and trusted-store assumptions. It does not guarantee semantic truth or runtime success.
  \item[What is trusted?] The coordinator code, durable store, locally configured trust policy, admitted registry descriptions, and contract validators. Malicious admitted agents remain outside scope.
  \item[Is fencing consensus?] Fencing alone is not consensus. The legacy PostgreSQL mode serializes through one database; the etcd mode separately uses Raft consensus to order authoritative state. Fencing then prevents stale session owners from committing and must also be enforced by effectful receivers.
  \item[Does the cluster scale automatically?] Membership of already running coordinator nodes is discovered and reconciled automatically. The software does not provision compute, and permanent majority loss requires operator-led disaster recovery.
  \item[Is there no single point of failure?] With at least three voters placed in independent failure domains, no single coordinator or etcd-member crash stops the authoritative path. One voter, one Docker host, shared network infrastructure, and a cluster-wide secret remain single-failure or correlated-failure risks.
  \item[How broad is A2A interoperability?] Controlled fixtures plus one vendored official sample. This is evidence of concrete compatibility, not universal conformance.
  \item[Did the hybrid work empirically?] Version 0.3 universally refused feasible observations. The constrained v0.4 bridge reached 90.0\% accuracy, 76.0\% feasible recall, zero false acceptance, and majority acceptance for all feasible cases. It still failed the pre-registered per-model floor because Qwen3 1.7B recall was 28\%.
  \item[Why trust the comparison?] The corpus and author labels were frozen before collection, labels were hidden during inference, all 540 identities were verified, and analysis is deterministic. However, labels were not independently adjudicated, so claims remain limited to framework conformance on this designed corpus.
\end{description}
